\documentclass[11pt]{article}
\usepackage{fontspec,amsmath,amssymb,amsthm,hyperref}
\usepackage[margin=1in]{geometry}
\setlength{\emergencystretch}{3em}
\title{The scalar SGLD contraction with the printed constants}
\author{LangevinLean mathematical exposition}
\date{}
\begin{document}\maketitle
\noindent Source module: \texttt{AverageLMC/SGLDRecurrence.lean}. Compile this proof document with XeLaTeX or LuaLaTeX. The accompanying HTML pairs the mathematical argument with the exact Lean statements and records the scope of the formalization.


Here m is the full curvature, R=n times the component RMS seminorm, and A is the full coordinate average. These scalar lemmas are applied to the actual-law recurrence in SGLDTheorem.

\[3R^2h\le m,\quad m^2\le R^2\quad\Longrightarrow\quad \frac{1-mh}{1+mh}+2R^2h^2\le\frac1{1+mh}.\]
\section{Compare the exact contraction and bias coefficients}


The restriction 3 R squared h at most m, together with m squared at most R squared, implies mh at most one third. Multiplying by the positive denominator 1+mh bounds the actual contraction by its reciprocal. The same step restriction bounds the persistent-error coefficient by the precise expression A+2m+2 R squared/m, divided by 1+mh.



\noindent Lean declarations: \texttt{sgld\_dimensionless\_step}, \texttt{sgld\_contraction\_le}, \texttt{sgld\_bias\_le}.

\section{Iterate around the exact fixed floor}


Set q=(1+mh) inverse and c=(A+2m+2 R squared/m)hp/m. The previous comparisons put the actual recursion below q times its preceding value plus (1-q)c. The finite-horizon affine induction gives q to the kth power times the initial value plus c, including the zero-step case.



\noindent Lean declarations: \texttt{sgld\_recurrence\_bound}.

\end{document}
